\documentclass[11pt]{article}
\usepackage[a4paper,margin=1in]{geometry}
\usepackage{xcolor}
\usepackage{enumitem}
\usepackage{hyperref}
\hypersetup{colorlinks=true,linkcolor=blue,urlcolor=blue}

\newcommand{\dSABRE}{d\textsc{sabre}}
\newcommand{\TeleSABRE}{Tele\textsc{sabre}}

% Reviewer comment / author response formatting
\newenvironment{comment}{\medskip\noindent\itshape}{\par}
\newcommand{\response}{\smallskip\noindent\textbf{Response.}\ }
\newcommand{\where}[1]{\textcolor{blue}{[#1]}}

\title{Response to Reviewers\\[2pt]
\large Manuscript TCAD-2026-0901\\
``\dSABRE{}: A SABRE-Style Router for Multi-Core Distributed Quantum
Computers''}
\author{Sanjiang Li}
\date{}

\begin{document}
\maketitle

\noindent
We thank the three reviewers and the editors for a careful and
constructive set of reviews.  The revision responds to all points; the
substantive changes are marked in blue in the manuscript, and
\where{like this} in this letter to point at where each change lands.

\section*{Summary of changes}

\begin{enumerate}[leftmargin=*]

\item \textbf{A unifying design principle, stated first and proved.}
The reviewers' central shared concern was that \dSABRE{} reads as a
collection of heuristics rather than a principled formulation.
Section~III now opens with a single routing potential $\Phi$, and
Proposition~1 proves that both the intra-core SWAP score and the
inter-core teleport score are the same quantity---the change in $\Phi$,
charged for resources spent and penalised for capacity consumed.  Every
scoring term is now derived from that principle rather than introduced
as a separate mechanism.

\item \textbf{Three propositions replace purely empirical
justification.}  Proposition~1 (unified objective), Proposition~2 (the
BFS-layer extended set is the ordering that makes a depth discount
well-defined, which topological filling cannot provide), and
Proposition~3 (forced progress, bounding a stall at twice the
core-graph diameter in EPR pairs under a stated capacity condition,
and separating unconditional termination from completion).  Each is attached to an
experimental result it explains rather than standing alone.

\item \textbf{A negative result we report rather than hide.}  We tested
whether the hop-gain term $g_\mathrm{hop}$ could be retired as
redundant, since the existing ablation showed it inert.  Re-running
every suite with $w_h{=}0$ showed the ablation subset was
misleading---it excludes the three largest circuits---and that removing
the term costs $18.6\%$ EPR on the 13\,040-CX multiplier and $2.2\%$ on
the 64-qubit geometric mean.  The term is kept, and Section~III-D now
justifies it by measurement and explains the scale dependence, which we
believe answers the reviewers' question about term necessity better
than a deletion would have.

\item \textbf{Mechanism roles reported honestly.}  The capacity penalty
is now presented as the central mechanism it measurably is; congestion
relief and checkpoint--rollback are reframed as \emph{feasibility}
mechanisms and reported in terms of completion rather than EPR
reduction, since on the corrected initial layout their effect on the
geometric mean is neutral.

\item \textbf{A non-grid architecture.}  A new experiment runs the
64-qubit circuits on a ring of four IBM 27-qubit heavy-hex cores, with
no router or layout change and no parameter retuned.  \dSABRE{}'s
margin over \TeleSABRE{} is $48.2\%$ there, slightly larger than on the
grid architectures, and \TeleSABRE{} additionally fails to converge on
a circuit \dSABRE{} routes.

\item \textbf{pytket-dqc rebuilt to our physical architecture.}  We
found that the network description \texttt{pytket-dqc} is normally
given is more permissive than the device \dSABRE{} is benchmarked on,
in two correctable respects.  Its per-server communication capacity
defaults to \emph{unbounded}, whereas a core has one communication port
per inter-core link ($2$--$3$ here); and the server sizes passed are an
even split of the logical qubit count rather than the physical capacity
$M-\deg(c)$.  Section~IV-C re-runs it with both corrected, and Table~IV now reports
the corrected figures throughout.  On the smaller suites the data
correction \emph{helps} \texttt{pytket-dqc}: its geometric mean improves
from $35.3$ to $27.5$ at 25 qubits and $12.3$ to $10.6$ at 36, so
\dSABRE{}'s margin narrows to $-44.3\%$ and the 36-qubit comparison
crosses to parity ($+1.8\%$).  At 64 qubits it goes the other way
($168.0$ to $176.3$), narrowing \dSABRE{}'s deficit from $+20.0\%$ to
$+14.7\%$.  We report all of these in place of the earlier numbers.

Separately, the communication bound is not enforced by the distributors
at all: setting it leaves every reported cost unchanged, because the
check happens only at circuit materialisation.  Auditing that,
\textbf{13 of the 21 distributions cannot be built within the
communication capacity the device provides}---AE alone needs at least
$19$ link qubits per core against the $2$ available.  We think this is
the most useful single fact we can give a reader weighing the two
columns, and it is why the e-bit comparison is no longer a headline
figure.

\item \textbf{Comparison methodology, and a change of emphasis.}  A new
table states what one EPR pair buys under each compiler's cost model,
and the differing primitives are addressed by a layout-controlled
comparison and an entanglement-lifetime sweep.  Following the
reviewers' concerns we also changed which comparison the paper
headlines.  \TeleSABRE{} shares \dSABRE{}'s cost model, physical
architecture and intra-core SWAP accounting, so it is now the only
cross-compiler figure quoted in the abstract, introduction and
conclusion.  \texttt{pytket-dqc}'s e-bit counts are reported in full in
Section~IV together with the four respects in which they are not equal
work---no intra-core routing, insensitivity to interconnect bandwidth,
unbounded entanglement lifetime, and gate grouping---but no longer as a
headline number.  We removed the favourable \texttt{pytket-dqc}
figures from the abstract along with the unfavourable ones, since both
rest on the same uncontrolled cost model.

\item \textbf{Presentation.}  New notation table; rewritten abstract;
reorganised introduction; consistent \dSABRE{} casing throughout;
abbreviations expanded at first use; Figure~4 caption rewritten;
compile-time claim justified.

\end{enumerate}

\subsection*{An important note on the reported numbers}

While preparing this revision we found and fixed a defect in the
initial-layout stage that affected the numbers in the submitted
manuscript.  The corner-reservation rule was written for a monolithic
chip: it removed the four globally lowest-degree physical qubits, which
under core-major qubit numbering all lie in core~0, so every other core
was left with no reserved communication or escape slot and was
regularly packed full by \texttt{SabreLayout}.  The fix reserves the
most-remote corners of \emph{every} core, with the count chosen
adaptively from the fill ratio.

All benchmark tables have been regenerated on the corrected layout.
The correction \emph{strengthens} the headline results and removes two
weaknesses a reader could reasonably have questioned: the GHZ-64q
regression is now a win, and \dSABRE{} now beats \TeleSABRE{} on the
100-qubit QFT it previously lost.  We also upgraded the
\texttt{pytket-dqc} baseline from its partitioning allocator to its
strongest Steiner-embedding distributor; this is a harder baseline that
wins several structured circuits under an unbounded entanglement
lifetime, and Section~IV reports that outcome plainly rather than
selecting the weaker configuration.  We judged that presenting the
strongest available baseline was the right response to the reviewers'
questions about comparison fairness, even though it narrows some
margins.

%%─────────────────────────────────────────────────────────────────────
\section*{Reviewer 1}

\begin{comment}
The overall design of dSABRE appears to be primarily driven by a
collection of heuristic optimizations addressing several observed
failure modes rather than being derived from a unified optimization
principle\ldots{} the method reads more like a set of incremental fixes
layered onto a baseline routing framework.
\end{comment}

\response
We agree this was the manuscript's main weakness, and it is the change
we invested most in.  Section~III-A now states the design principle
before any mechanism appears: \dSABRE{} chooses the action minimising
$\mathrm{cost}(a) + \mathrm{pen}(a) + \Delta\Phi(a)$, where $\Phi$ is a
single circuit-wide potential combining immediate front-layer distance
with depth-discounted lookahead.  Proposition~1 proves that both
scoring rules in the paper are instances of this expression.
\where{Section~III-A, Eqs.~(2)--(3), Proposition~1}

Crucially, the revision also draws a line that was previously blurred.
Congestion relief and checkpoint--rollback are \emph{not} additional
score terms; they exist because a greedy minimiser can leave the region
where $\Phi$ is a usable surrogate.  We now say so explicitly and
report their value as feasibility rather than EPR reduction.
\where{end of Section~III-A; Contributions bullet~3}

\begin{comment}
The exclusive use of teledata, while TELESABRE supports both teledata
and telegate, is not sufficiently motivated\ldots{} raises questions
about fairness and completeness of the comparison.
\end{comment}

\response
Three points, all now in the manuscript.  First, on fairness: the
restriction handicaps \dSABRE{}, not \TeleSABRE{}.  \dSABRE{} uses a
strict subset of \TeleSABRE{}'s primitives and still wins on the large
majority of circuits, so the comparison understates rather than
overstates the result.  Second, on motivation: teledata is not merely a
simplification but the lifetime-robust primitive.  A teledata EPR pair
is consumed instantaneously by its teleport, whereas gate teleportation
requires a link to remain entangled for the span of the gate group it
serves; the new lifetime sweep quantifies exactly what that assumption
is worth.  Third, on completeness: we agree telegate support is a real
extension, and it is now stated as scoped future work rather than left
implicit.
\where{Section~II-B; Section~IV-A ``What one EPR pair buys'' and
Table~III; lifetime sweep in Section~IV-B; Section~VI}

\begin{comment}
The architecture is described as an abstract graph GA treated as a
black box, yet the experimental setup relies specifically on a grid-of-grids
structure\ldots{} the reported performance differences across circuit
families suggest that the method is sensitive to underlying topology.
\end{comment}

\response
This was a fair criticism of the evidence, and we have added the
missing experiment rather than argued the point.  Section~IV now
reports the 64-qubit circuits on a ring of four IBM 27-qubit heavy-hex
cores---irregular degree, six degree-1 leaves, intra-core diameter 12
against 6 for the grid.  No router or layout code is specialised for
it: the corner-reservation rule is defined by vertex degree and
remoteness, so it selects leaf qubits instead of grid corners
automatically, and every hyperparameter keeps its default.
\TeleSABRE{} runs on the same architecture.  \dSABRE{}'s margin is
$48.2\%$, slightly \emph{larger} than on the grids.

The absolute numbers also behave in an interpretable way that we now
report: EPR counts fall for both routers while intra-core SWAP counts
rise $50$--$70\%$, because four 27-qubit cores partition the same 64
logical qubits more coarsely than six 16-qubit cores but are twice as
wide.  Larger cores trade inter-core EPR pairs for intra-core SWAPs,
and \dSABRE{} makes that trade without retuning because both are scored
against the same potential.
\where{new Section~IV-C, Table~V}

On topology sensitivity specifically: the GHZ regression the reviewer
noted was in fact a symptom of the initial-layout defect described
above, not of the scoring rule.  With every core given reserved escape
capacity, GHZ-64q improves rather than regresses.  We report this
mechanism explicitly.
\where{Section~IV-B, 64q discussion}

\begin{comment}
The mathematical formulation does not fully clarify why this particular
decomposition of terms is necessary, nor how each term contributes to a
unified optimization goal.
\end{comment}

\response
Addressed by the same restructure, and by testing the question
empirically rather than asserting an answer.  Four of the five terms
are now \emph{derived}: Proposition~1 shows they are exactly
$\mathrm{cost}$, $\mathrm{pen}$, and $\Delta\Phi$.  The fifth,
$g_\mathrm{hop}$, is the one term the proposition does not imply, so we
tried to remove it.  That experiment failed in an informative way: the
term is exactly inert on all twelve B-grid circuits, but removing it
costs $18.6\%$ EPR on the largest circuit in the suite and $2.2\%$ on
the 64-qubit geometric mean, because $\Delta_F$ is evaluated at a
specific landing port and under-reports core-level progress as
intra-core distances grow.  We keep the term and now say exactly what
it buys and where.  We also note the methodological trap, since it bears
on how ablations should be read: measured only on the six-circuit
subset the term looks like noise.
\where{Section~III-D; Section~IV-D and Table~VI}

\begin{comment}
A more explicit discussion of the equivalence (or lack thereof) between
different EPR usage models would help strengthen the validity of the
reported improvements.
\end{comment}

\response
Agreed; this discussion was scattered across the experimental section
and two appendices.  It is now consolidated into one place with a
table stating, for each compiler, what a single EPR pair purchases and
whether intra-core SWAP cost is modelled at all.  We then address the
non-equivalence in two ways rather than only flagging it: a
layout-controlled comparison in which \dSABRE{} and the baseline start
from the same partition, and a sweep bounding how many gates one EPR
pair may serve.
\where{Section~IV-A and Table~III; Section~IV-C and Table~V
(physical-architecture comparison); layout-controlled results and
lifetime sweep in Section~IV-B}

%%─────────────────────────────────────────────────────────────────────
\section*{Reviewer 2}

We thank the reviewer for the detailed and balanced assessment, and in
particular for the observation about relative mechanism importance,
which prompted a substantive change in how the paper reports its own
results.

\begin{comment}
The level of novelty appears moderate\ldots{} many of the contributions
can be viewed as heuristic refinements of existing routing
methodologies rather than fundamentally new algorithmic concepts.
\end{comment}

\response
We have repositioned the claim rather than inflating it.  The
contribution is not a longer score---after this revision it is a
shorter one---but the observation that a distributed router's two
movement decisions can and should be derived from one potential.  Prior
SABRE-style distributed routers design the local SWAP heuristic and the
inter-core energy as separate objects, with no guarantee they agree
about which qubit should move; unifying them is what allows the
architecture's capacity constraint to enter the score at all, and the
ablation attributes more of the improvement to that single signal than
to anything else.  Related Work now argues this distinction directly
instead of enumerating feature differences.
\where{Contributions bullet~1; Section~V, ``SABRE-style distributed
routers''}

\begin{comment}
The paper relies predominantly on empirical evidence\ldots{} there is
limited theoretical insight into why the proposed scoring formulation
should consistently outperform alternative approaches.
\end{comment}

\response
The revision adds three propositions, each tied to a specific
experimental result so that theory and evidence explain one another
rather than sitting in separate sections.  Proposition~1 (unification)
is what licenses the removal of a redundant term.  Proposition~2 shows
BFS layering is the \emph{only} construction under which the depth
discount is well-defined, with a counterexample family showing
topological filling can exclude a gate that becomes executable one layer
ahead---this predicts the ablation result and its growth with circuit
size.  Proposition~3 bounds the cost of escaping a stall at twice the
core-graph diameter in EPR pairs, and is deliberately careful about
what it does and does not promise: termination is unconditional, but
completion rests on a capacity condition that can fail if a region of
the chip saturates.  This is what explains \dSABRE{} completing
circuits on which a fixed safety-valve cap fails, without claiming more
than the algorithm delivers.
\where{Propositions~1--3 in Sections~III-A, III-E, III-G; proofs of 2 and 3 in the appendices}

\begin{comment}
Since the considered approaches employ different communication
primitives, direct comparisons based solely on EPR counts do not always
provide a fully controlled evaluation.
\end{comment}

\response
See the response to Reviewer~1's fifth point.  Beyond the consolidated
cost-model discussion, we added two controlled comparisons: identical
initial partitions for both compilers, and a bound on gates served per
EPR pair.  The second is the more informative: under an unbounded
lifetime the gate-teleportation baseline wins several structured
circuits, but its advantage crosses parity between 10 and 5 gates per
pair and inverts at the value proposed in the recent literature.  We
report both regimes.
\where{Section~IV-B, layout-controlled paragraph and the
lifetime sweep}

\begin{comment}
The ablation study indicates that the capacity-aware penalty
contributes a significant fraction of the observed gains, whereas some
of the remaining mechanisms provide comparatively modest improvements.
A clearer discussion of the relative importance of the different design
choices would help readers.
\end{comment}

\response
This observation was correct and we have restructured the paper's
claims around it.  The capacity penalty is now presented as the central
mechanism throughout---abstract, contributions, and ablation---rather
than as one of three co-equal ideas.  The hop-gain term, being inert,
is removed entirely.  Congestion relief, which the submitted version
credited with a large geometric-mean effect, is on the corrected
initial layout neutral; we now say so plainly and reclassify it, with
checkpoint--rollback, as a feasibility mechanism whose value is
completing circuits that otherwise fail.  We would rather report this
than retain an overstated claim.
\where{Abstract; Contributions bullets~1 and~3; Section~IV-D and
Table~VI}

\begin{comment}
Incorporating latency, scheduling effects, fidelity considerations, or
more application-oriented distributed workloads would further
strengthen the practical relevance.
\end{comment}

\response
Partly addressed, partly scoped.  Three application-oriented circuits
(QPE, QAOA, and a 13\,040-CX arithmetic multiplier) are added to the
64-qubit suite, taking the evaluation to 21 circuits.  Fidelity enters
through the lifetime bound, which is a fidelity-motivated constraint:
capping gates-per-EPR is how a decaying link appears in a cost model.
Latency is partly visible already, since intra-core SWAP count---the
dominant contributor to routed depth---is reported beside EPR count in
every table and falls alongside it.  A genuine joint objective over
EPR, depth, and error rates would require a device error model and a
scheduling pass inside the loop, both outside the routing stage this
paper addresses; we now state that limitation explicitly rather than
leaving it implicit.
\where{Table~IV (64q panel); Section~IV-B; new paragraph in Section~VI}

%%─────────────────────────────────────────────────────────────────────
\section*{Reviewer 3}

We thank the reviewer for the detailed presentation review; these
comments led to the most visible changes in the manuscript.

\begin{comment}
1. The Abstract should be rewritten to clearly state the motivation for
the problem and provide a concise summary of the main contributions.
\end{comment}

\response
The abstract is rewritten in three paragraphs: what a DQC is and why
EPR consumption is the objective; the design principle and the three
consequences that follow from it; and the results.  Ablation
percentages have been removed from the abstract, and all abbreviations
are expanded.
\where{Abstract}

\begin{comment}
2. The overall organization needs improvement\ldots{} after introducing
the second problem class, the discussion becomes overly detailed,
causing the reader to lose track of the overall structure before the
third class is presented.
\end{comment}

\response
The three problem classes are now presented in parallel form, one or
two sentences each, with no digression between them.  The material on
communication primitives that previously interrupted class~(ii) has
moved to Background~II-B, where those primitives are properly
introduced, and the paragraph on why EPR consumption is the objective
now follows the enumeration instead of being embedded inside it.
\where{Section~I, paragraphs 2--3}

\begin{comment}
3. In the Contributions section, the notations $\Delta_F$ and
$\Delta_E$ are introduced without definition\ldots{} the term `comm
port' should be defined or expanded upon.
\end{comment}

\response
The contributions are now written in prose with no undefined symbols;
$\Delta_F$ and $\Delta_E$ first appear in Section~III where they are
defined.  ``Comm port'' is expanded to ``communication port''
throughout.  More broadly, and following the reviewer's overall
recommendation, all notation is now collected in a single table at the
start of Section~III.
\where{Contributions; Table~I}

\begin{comment}
4. In the Background section, the statement ``Together, CX and
arbitrary 1Q gates form a universal gate set'' should specify the 1Q
gate sets.
\end{comment}

\response
Corrected and made precise: CX together with the full single-qubit
unitary group $U(2)$ is exactly universal, with a citation; a finite
basis such as $\{\mathrm{CX},H,T\}$ suffices for approximate
universality, and we note this is the basis the benchmark circuits are
transpiled into.
\where{Section~II-A}

\begin{comment}
5. The claim that DSABRE is slower than TeleSABRE for circuit
compilation requires further justification.
\end{comment}

\response
Expanded into a justified argument with three strands: the measured
growth rate matches the derived complexity rather than exceeding it, so
the gap is a bounded constant across a $200\times$ range of circuit
size rather than an asymptotic penalty; the two implementations differ
exactly as that constant predicts (compiled C\texttt{++} with cached
distance lookups versus pure Python calling NetworkX inside the routing
loop); and the ordering inverts on Graphstate, where \dSABRE{} is
$78\times$ faster, showing the constant is not intrinsic.  We also note
that nothing in the complexity analysis is language-specific, and that
the absolute cost is modest against the EPR savings it buys.
\where{Section~IV-F}

\begin{comment}
6. The caption of Figure 4 is unclear and should be revised.
\end{comment}

\response
Rewritten.  The caption now opens with one sentence saying what the
figure shows and what decision is being made, then describes the two
execution steps, and only then walks the three candidates---each of
which now isolates a single decision factor.  The dense parenthetical
definitions have moved into the running-example text.
\where{Figure~4 caption; Table~II}

\begin{comment}
Editorial 1. The acronym DSABRE is written inconsistently\ldots{}
Editorial 2. Abbreviations should be expanded when they first appear.
\end{comment}

\response
Both fixed.  The name is now set by a single macro rendering
``\dSABRE{}'' with a lowercase initial everywhere, matching the title;
the same was done for \TeleSABRE{}.  EPR, DAG, MQT-Bench, DQC, QPU and
the remaining abbreviations are expanded at first use, including in the
abstract.
\where{Throughout}

\begin{comment}
I strongly recommend defining all notation systematically in one place.
\end{comment}

\response
Done: Table~II collects every symbol used in the technical section with
its meaning and defining equation.
\where{Table~I}

%%─────────────────────────────────────────────────────────────────────
\section*{Area Editor and Associate Editor}

The four required revisions map onto the changes above:
sharpening the core novelty (Summary items~1, 2 and Reviewer~2's first
point); justifying design choices and how they cohere (items~1--4);
improving exposition and notation (item~7 and all of Reviewer~3);
and strengthening the experimental comparisons with attention to
differing communication primitives and EPR cost models (items~5, 6).

We note in closing that this revision removes a scoring term, reports a
mechanism as neutral that the submitted version credited with a large
effect, and adopts a stronger baseline that wins several circuits.  We
believe a paper that states these plainly is more useful than one that
does not, and that the central result---a unified scoring principle in
which architectural capacity is a first-class term---is better
supported for it.

\end{document}
